\subsection{Docker Compose i konfiguracja silników}

Środowisko testowe zostało sprowadzone do jednego pliku \texttt{docker-compose.yml}
definiującego cztery izolowane kontenery -- po jednym na każdy z porównywanych silników
bazodanowych. Pozwoliło to uzyskać pełną reprodukowalność (każde uruchomienie
startuje z identycznego obrazu), czyste odseparowanie procesów (żaden silnik nie
dzieli RAM-u innym) oraz jednolity interfejs startu i zatrzymania
całego stosu jedną komendą \texttt{docker compose up}.

Każdy kontener ma jawnie ustawiony limit pamięci (\texttt{deploy.resources.limits.memory})
oraz własne, świadomie dobrane opcje tuningowe -- tak aby zapewnić porównywalne
warunki pracy bez nieuczciwego faworyzowania żadnej z baz.
Wszystkie kontenery udostępniają swoje porty
(\texttt{5432}, \texttt{3306}, \texttt{27017}, \texttt{7687}+\texttt{7474}),
dzięki czemu warstwa orkiestracji w Pythonie łączy się do nich przez \texttt{localhost}.
Dane są utrzymywane  w czterech nazwanych wolumenach Docker (\texttt{postgres\_data},
\texttt{mysql\_data}, \texttt{mongo\_data}, \texttt{neo4j\_data}), co pozwala zachować
załadowaną bazę między kolejnymi uruchomieniami benchmarków.

\subsubsection*{PostgreSQL 17}

Kontener \texttt{ztbd-postgres} oparty o obraz \texttt{postgres:17} z limitem 2~GB
RAM. Schemat bazy (12 tabel wraz z kluczami obcymi i kaskadami jest ładowany automatycznie przy pierwszym starcie
kontenera, poprzez zamontowanie pliku \texttt{init.sql} w katalogu
\texttt{/docker-entrypoint-initdb.d/}. Parametry uruchomieniowe silnika:

\begin{lstlisting}[language=bash,basicstyle=\ttfamily\small]
shared_buffers=512MB
work_mem=16MB
maintenance_work_mem=512MB
effective_cache_size=1536MB
wal_buffers=16MB
checkpoint_completion_target=0.9
max_wal_size=2GB
synchronous_commit=off
\end{lstlisting}

\subsubsection*{MySQL 8.0}
Kontener \texttt{ztbd-mysql} oparty o obraz \texttt{mysql:8.0}, również z limitem
2~GB RAM i własnym \texttt{init.sql}. Parametry silnika InnoDB:

\begin{lstlisting}[language=bash,basicstyle=\ttfamily\small]
--innodb-buffer-pool-size=1073741824    (1 GB)
--innodb-log-file-size=256M
--innodb-flush-log-at-trx-commit=2
--innodb-flush-method=O_DIRECT
--innodb-doublewrite=0
--innodb-change-buffering=all
--bulk-insert-buffer-size=256M
--local-infile=1
--max-allowed-packet=256M
--innodb-autoinc-lock-mode=2
\end{lstlisting}


\subsubsection*{MongoDB 8.0}

Kontener \texttt{ztbd-mongo} oparty o obraz \texttt{mongo:8.0}, limit 2~GB RAM,
port \texttt{27017}. Konfiguracja silnika WiredTiger sprowadza się do jednej,
ale kluczowej opcji:

\begin{lstlisting}[language=bash,basicstyle=\ttfamily\small]
mongod --wiredTigerCacheSizeGB 1.2
\end{lstlisting}

\noindent 1.2~GB to 60\% limitu kontenera -- rozmiar cache silnika pamięciowego
WiredTiger, w którym MongoDB trzyma skompresowane strony dokumentów i indeksów.
Pozostałe 40\% pozostawione jest na procesy Mongo poza cache (połączenia, sortowanie,
agregacje, heap). Schemat nie jest ładowany -- Mongo nie wymaga definicji kolekcji
z góry, tworzą się one przy pierwszym zapisie.

\subsubsection*{Neo4j 2026.02.2}

Kontener \texttt{ztbd-neo4j} oparty o obraz \texttt{neo4j:2026.02.2} (najnowsza
wersja z aktywną linią wsparcia) z doinstalowanym pluginem \textbf{APOC}.
Ze wszystkich
czterech baz Neo4j dostał \textbf{najwyższy limit pamięci -- 4~GB} -- ze względu
na architekturę opartą o JVM, która wymaga osobnego budżetu na heap, osobnego
na page cache plików bazy oraz osobnego na pamięć transakcyjną:

\begin{lstlisting}[language=bash,basicstyle=\ttfamily\small]
server.memory.heap.initial_size      = 1g
server.memory.heap.max_size          = 1g
server.memory.pagecache.size         = 1g
dbms.memory.transaction.total.max    = 3g
\end{lstlisting}


% -------------------------------------------------------
\subsection{Technologie i biblioteki}

Warstwa orkiestracji (generowanie danych, ładowanie, benchmarki, \texttt{EXPLAIN},
wizualizacja) została zaimplementowana w \textbf{Pythonie 3.12+}. Wybór Pythona
wynika z dostępności natywnych, dojrzałych driverów do wszystkich czterech silników,
bogatego ekosystemu narzędzi analitycznych (\texttt{pandas}, \texttt{matplotlib})
oraz niskiego narzutu zapisu kodu -- zestaw 24~scenariuszy da się zwięźle wyrazić
w języku wysokiego poziomu, bez zaciemniania logiki pomiaru ogólnikami języka.

Cały projekt jest uruchamiany przez jeden punkt wejścia (\texttt{main.py}) udostępniający
sześć komend CLI zbudowanych na module \texttt{argparse}: \texttt{generate},
\texttt{load}, \texttt{benchmark}, \texttt{explain}, \texttt{visualize} oraz
\texttt{run-all} (pełny pipeline 8-krokowy: generowanie~$\rightarrow$~ładowanie
bez indeksów~$\rightarrow$~benchmarki~$\rightarrow$~\texttt{EXPLAIN}~$\rightarrow$~ładowanie
z indeksami~$\rightarrow$~benchmarki z indeksami~$\rightarrow$~\texttt{EXPLAIN} z
indeksami~$\rightarrow$~wizualizacja).

\subsubsection*{Drivery bazodanowe}

\begin{itemize}
    \item \textbf{psycopg 3} -- PostgreSQL. Najnowszy driver, obsługujący operacje asynchroniczne  
    protokół \texttt{COPY FROM STDIN} strumieniowo,
    co jest kluczowe dla wydajnego ładowania danych.
    \item \textbf{PyMySQL} -- MySQL. Driver napisany w czystym Pythonie (bez zależności na
    \texttt{mysqlclient} wymagającym kompilacji); wspiera \texttt{LOAD DATA LOCAL INFILE}
    poprzez flagę \texttt{local\_infile=True} w \texttt{connect()}.
    \item \textbf{PyMongo} -- MongoDB. Oficjalny driver do MongoDB; wspiera
    \texttt{insert\_many(ordered=False)} (kontynuacja po błędach duplikatów) oraz
    korzystanie z zaawansowanych zapytań typu \texttt{aggregate}.
    \item \textbf{neo4j} (oficjalny driver Pythonowy) -- Neo4j. Komunikacja przez
    protokół Bolt na porcie 7687; kluczowe jest tu wykorzystywanie jednej sesji
    (\texttt{driver.session()}) dla wielu operacji -- żeby uniknąć narzutu połączeń TCP przy każdym batchu.
\end{itemize}

\subsubsection*{Generowanie danych}

Generator danych (\texttt{src/generators/data\_generator.py}) opiera się na
bibliotece \textbf{Faker} z \textbf{ustalonym seedem 42}.
Stały seed gwarantuje pełną powtarzalność -- każde wywołanie
\texttt{python main.py generate --volume small} produkuje identyczny
zestaw plików CSV, co jest niezbędne dla rzetelności porównania (wszystkie cztery bazy
dostają dokładnie te same dane wejściowe). Wyjściem są pliki CSV
(po jednym na tabelę), które każdy z czterech loaderów czyta niezależnie
i transformuje do swojej natywnej reprezentacji.

\subsubsection*{Analiza i wizualizacja}

Warstwa analityczna korzysta z trzech standardowych bibliotek:

\begin{itemize}
    \item \textbf{pandas} -- wczytywanie plików CSV z wynikami, agregacje
    (\texttt{groupby} po scenariuszu, bazie, wariancie indeksów), obliczanie
    mediany, min i max z 3 triali.
    \item \textbf{NumPy} -- operacje macierzowe przy liczeniu heatmap.
    \item \textbf{matplotlib} (backend \texttt{Agg}) -- generowanie wykresów PNG
    bez wymagania GUI.
\end{itemize}


